%! The Tangent Function — Unit 7, Lesson 7.1 — Exit Ticket (Answer Key)
\documentclass[10pt]{article}
\usepackage{precalculus-article}
\usepackage{precalculus-key}   % pulls in -boxes; do NOT also load -boxes
\newcommand{\TallMath}[1]{$\displaystyle #1\rule[-1.4em]{0pt}{3.2em}$}

\begin{document}

\pageheader{Unit 7, Lesson 7.1}{Exit Ticket --- Answer Key}
\namedateperiod

\vspace{0.10in}

\begin{notesbox}{Exit Ticket --- Answer Key}

\textbf{1.}\quad Given $\sin\theta = \dfrac{\sqrt{3}}{2}$ and $\cos\theta = \dfrac{1}{2}$,
compute $\tan\theta$.  Show your work.

\vspace{4pt}
$\tan\theta = \dfrac{\sin\theta}{\cos\theta} = $ \ans{$\dfrac{\sqrt{3}/2}{1/2} = \sqrt{3}$}

\vspace{0.14in}
\textbf{2.}\quad State the period and the location of the two positive vertical
asymptotes closest to $\theta = 0$ for the function $g(\theta) = \tan(2\theta)$.

\vspace{4pt}
Period of $g$: \ans{$\pi/2$}

\vspace{4pt}
Asymptotes: $\theta = $ \ans{$\pi/4$} and $\theta = $ \ans{$3\pi/4$}

\begin{teachernote}
Asymptotes of $\tan(2\theta)$: set $2\theta = \pi/2 + k\pi \Rightarrow \theta = \pi/4 + k\pi/2$.
The two smallest positive values are $\pi/4$ ($k=0$) and $3\pi/4$ ($k=1$).
\end{teachernote}

\vspace{0.14in}
\textbf{3.}\quad For the parent function $y = \tan\theta$, the graph is always
\ans{increasing} between consecutive
asymptotes, and changes from concave \ans{down} to concave
\ans{up} within each period.

\end{notesbox}

\end{document}
